\clearpage
\VerdierPhysicalPageAnchor{0187}%
\begingroup\small\setlength{\parskip}{0pt}\setlength{\emergencystretch}{1.5em}
\setlength{\abovedisplayskip}{4pt plus 2pt minus 2pt}%
\setlength{\belowdisplayskip}{4pt plus 2pt minus 2pt}%

\VerdierSemanticAnchor{ch3d-2-1-1-base-change}
Soit alors $m\colon Z\longrightarrow X$ un morphisme de complexes. Il est
clair que le complexe $Z'$ au-dessus de $Z$ obtenu par le changement de
base $Z\xrightarrow{m}X$ :

\VerdierSemanticAnchor{ch3d-2-1-1-base-change-square}
\begin{center}
\begin{tikzcd}[column sep=4.2em,row sep=2.6em]
Z' \arrow[r] \arrow[d] & X(u,i) \arrow[d,"{p(u,i)}"] \\
Z \arrow[r,"m"'] & X
\end{tikzcd}
\end{center}
est isomorphe à un complexe
$p(w,i)\colon Z(w,i)\longrightarrow Z$ pour un morphisme $w$ convenable.

\VerdierSemanticAnchor{ch3d-2-1-2}
\medskip\noindent\textbf{Lemme 2.1.2.}\quad
\textit{a) Le morphisme
$H^j(p(u,i))\colon H^j(X(u,i))\longrightarrow H^j(X)$ est un
isomorphisme pour $j\neq i,i+1$.}

\noindent\textit{b) On a des suites exactes :}

\VerdierSemanticAnchor{ch3d-2-1-2-exact-sequences}
\[
0\longrightarrow H^i(X(u,i))\longrightarrow H^i(X)
\longrightarrow
\frac{\operatorname{Ker}(d_X^i)+\operatorname{Im}(u)}
 {\operatorname{Im}(d_X^{i-1})+\operatorname{Im}(u)}
\longrightarrow0,
\]
\[
\begin{aligned}
0\longrightarrow{}&
\frac{\operatorname{Ker}(d_X^i)+\operatorname{Im}(u)}
 {\operatorname{Im}(d_X^{i-1})+\operatorname{Im}(u)}
\longrightarrow
\frac{X^i}{\operatorname{Im}(d_X^{i-1})+\operatorname{Im}(u)}\\
&\longrightarrow
\frac{X^i}{\operatorname{Ker}(d_X^i)+\operatorname{Im}(u)}
\longrightarrow0,
\end{aligned}
\]
\[
0\longrightarrow
\frac{X^i}{\operatorname{Ker}(d_X^i)+\operatorname{Im}(u)}
\longrightarrow H^{i+1}(X(u,i))
\longrightarrow H^{i+1}(X)\longrightarrow0.
\]

\noindent\textit{c) Pour que $p(u,i)$ soit un quasi-isomorphisme, il
faut et il suffit que :}
\[
\operatorname{Im}(d_X^{i-1})+\operatorname{Im}(u)=X^i.
\]

\VerdierSemanticAnchor{ch3d-2-1-2-proof}
La ``chasse au diagramme'' est laissée au lecteur.

\VerdierSemanticAnchor{ch3d-2-1-3}
\medskip\noindent\textbf{Définition 2.1.3.}\quad
Un morphisme de complexes $m\colon Z\longrightarrow X$ est appelé un
\emph{morphisme élémentaire de hauteur $i$} s'il est isomorphe dans
$\mathcal{F}\ell_1(\operatorname{comp}(\mathcal A))$ à un morphisme du type
$p(u,i)$, où $u$ est un épimorphisme.

\VerdierSemanticAnchor{ch3d-2-1-3-properties}
Un morphisme élémentaire est donc, en particulier, un quasi-isomorphisme
((2.1.2), c)). Un morphisme élémentaire est un épimorphisme de
$\operatorname{comp}(\mathcal A)$. Un épimorphisme de
$\operatorname{comp}(\mathcal A)$ est un morphisme élémentaire si et seulement
si son noyau est un complexe acyclique comportant au plus deux composants
non nuls. Le composé de deux morphismes élémentaires de même hauteur est
\endgroup
